\section{System Overview}
\label{sec:methodology-system-overview}
\begin{outline}
  Present a high-level overview of your proposed control framework,
  likely including a block diagram (similar to Figure 1 in the
  proposal) and a description of each component.
\end{outline}
This work presents a control framework for quadruped robots capable
of generating dynamic, acyclic gaits in challenging environments. The
framework employs a hierarchical architecture that processes robot
state and terrain data to produce footstep commands for the robot controller.

The first stage of the framework involves the user input
$\mathbf{v^{usr}}$ being supplied alongside the robot state to a gait
generation network (\textit{GaitNet}). GaitNet acts as an evaluative
surface, performing continuous optimization over the physical
workspace utilizing projected gradient ascent. Unlike discrete
sampling methods, GaitNet ranks these continuous outputs iteratively,
producing the optimal footstep action $\mathbf{f_a}$, which is passed
to the robot controller. The robot controller is the MIT Mini-Cheetah
Controller \cite{di_mini_cheetah_2020}, implemented
in python by \cite{zhuang_rl_mpc_locomotion_2025}.

This hierarchical design enables strict control over the range of
possible actions the robot can execute. Between the two stages,
candidate actions are filtered to ensure that only valid, prescribed
movements are permitted.

\begin{todo}
  Update \autoref{fig:diagram-control-system} with images of the
  experimental environment. Also update terminology for different
  control blocks and connections between them.
\end{todo}

\begin{figure}[H]
  \centering
  \includegraphics[width=1.0\linewidth]{images/diagrams/control-system.png}
  \caption{Block diagram of the proposed framework. Novel aspects
    shown in blue. The user defines an input direction
    $\mathbf{v^{usr}}$, which \textit{GaitNet}
    (\autoref{sec:methodology-gaitnet}) uses along with the robot state
    $\mathbf{x_g}$ and the current foot positions $\mathbf{p_f}$ to
    evaluate candidate footsteps. Instead of discrete sampling, GaitNet
    utilizes a continuous projected gradient ascent optimization
    strategy, generating the optimal footstep action $\mathbf{f_a}$
    to send to the
    robot controller. The \textit{MIT Mini-Cheetah Controller}
    \cite{di_mini_cheetah_2020} receives $\mathbf{f_a}$ and executes
  the corresponding actions.}
  \label{fig:diagram-control-system}
\end{figure}

\footnotetext{As is explained in later sections, GaitNet only outputs
  one footstep action at a time. The two separate foot target locations
  shown in the GaitNet visual would be selected sequentially with a
  1/25s delay. This visual is meant to highlight the near simultaneous
nature of the footstep selection process.}
